\documentclass[11pt]{article}
\usepackage[margin=2.0cm]{geometry}
\usepackage{booktabs}
\usepackage{amsmath}
\usepackage{amssymb}
\usepackage{enumitem}
\usepackage{array}
\usepackage{listings}
\usepackage{xcolor}
\usepackage[T1]{fontenc}
\usepackage{lmodern}

\setlength{\parindent}{0pt}
\setlength{\parskip}{3.6pt}
\newcommand{\code}[1]{\texttt{#1}}
\newcommand{\hr}{\vspace{2pt}\hrule\vspace{6pt}}
\renewcommand{\arraystretch}{1.12}

\definecolor{kw}{RGB}{0,90,160}
\definecolor{cmt}{RGB}{110,110,110}
\lstdefinestyle{vhdl}{
  language=VHDL, basicstyle=\ttfamily\footnotesize,
  keywordstyle=\color{kw}\bfseries, commentstyle=\color{cmt}\itshape,
  showstringspaces=false, columns=fullflexible, breaklines=true,
  frame=leftline, framesep=6pt, xleftmargin=8pt, aboveskip=6pt, belowskip=6pt
}
\lstset{style=vhdl}

\begin{document}

\begin{center}
{\Large\bfseries EE\,362 --- Computer Organization Laboratory}\\[2pt]
{\large\bfseries Lab 7 --- Guided Examples: Solution}\\[2pt]
{\itshape Ex 1 the program counter \quad$\cdot$\quad Ex 2 a real branch target \quad$\cdot$\quad Ex 3 two ways to get it wrong}\\[4pt]
Fall 2026 \quad$\cdot$\quad Lab TA: Gandhi Torres
\end{center}
\hr

\begin{center}
\small
\begin{tabular}{cl>{\raggedright\arraybackslash}p{6.0cm}>{\raggedright\arraybackslash}p{3.8cm}}
\toprule
& circuit & new idea & feeds the assignment as \\
\midrule
1 & \code{pc\_counter} & the PC is a \textbf{register with an adder in front} & \code{NIAG\_Reg} \\
2 & \code{branch\_target} & \textbf{PC+4, sign-extend, shift, add} & \code{NIAG\_NSA}, \code{NIAG\_Rel} \\
3 & \code{addressing\_traps} & off by a \textbf{factor} vs off by a \textbf{constant} & what your testbench must catch \\
\bottomrule
\end{tabular}
\end{center}

\vspace{4pt}\hr

\section*{Example 1 --- The program counter}

\begin{quote}
\textbf{A program counter is a register with an adder in front of it.} You built
the register in Lab 2; the adder is one line.
\end{quote}

\[
PC^{+} = \overline{\code{Rst}}\cdot\left( \code{Ld}\cdot \code{NextPC} + \overline{\code{Ld}}\cdot PC \right)
\]

That is Lab 2's equation with \code{Ld} renamed. \textbf{Only the reset value
changed.}

\subsection*{1.1 Why a register and not a counter}

A counter only counts. A program counter must also \textbf{jump} --- to a branch
target, to a subroutine, back again. So it is a \emph{loadable register}, and
``count'' is merely the special case of loading $PC+4$.

\begin{quote}
That is why Lab 8 replaces this \code{+4} with a multiplexer over five possible
next addresses. \textbf{Sequential execution is not a separate mechanism; it is
one input of the mux.}
\end{quote}

\subsection*{1.2 Why the reset vector matters}

On reset the PC must go somewhere \textbf{defined}. If it powers up at \code{'U'}
the CPU fetches garbage forever and every waveform you look at is meaningless.
Here, and in your CPU, the reset vector is address 0.

\textbf{In VHDL} the value is a \code{generic} used as a number, not a width:

\begin{lstlisting}
generic ( RESET_VEC : natural := 0 );
...
pc_reg <= std_logic_vector(to_unsigned(RESET_VEC, 32));
\end{lstlisting}

The example is 8 bits wide purely so the waveform is readable. Your
\code{NIAG\_Reg} is 32.

\vspace{4pt}\hr

\section*{Example 2 --- A real branch target}

This takes an actual instruction out of the ROM program your CPU runs in Lab 12.

\begin{center}
\begin{tabular}{lll}
\toprule
address & encoding & instruction \\
\midrule
\code{0x0C} & \code{0x18200002} & \code{BLEZ r1, +2} \\
\bottomrule
\end{tabular}
\end{center}

\subsection*{2.1 Field extraction}

\begin{center}
\begin{tabular}{lll}
\toprule
bits & field & value \\
\midrule
31--26 & opcode & \code{000110} $=$ \code{BLEZ} \\
25--21 & \code{rs} & \code{00001} $=$ \code{r1} \\
20--16 & \code{rt} & \code{00000} \\
15--0 & immediate & \code{0x0002} \\
\bottomrule
\end{tabular}
\end{center}

\subsection*{2.2 The computation}

\begin{align*}
\code{NSA}  &= PC + 4 = \code{0x0C} + 4 = \code{0x10} \\
\text{offset} &= \text{SignExt}(\code{0x0002}) \ll 2 = 2 \times 4 = \code{0x08} \\
\text{target} &= \code{NSA} + \text{offset} = \code{0x18} \qquad \text{--- the \code{SW} instruction}
\end{align*}

\subsection*{2.3 The two questions worth answering properly}

\textbf{Why is the base $PC+4$ and not $PC$?} Because by the time the branch is
evaluated the fetch has already happened --- the hardware has moved on. Using $PC$
would mean every target was one instruction earlier than the assembler wrote, so
the assembler would have to compensate. \textbf{Basing on $PC+4$ keeps the
encoding honest:} ``+2'' really means ``skip 2 instructions''.

\textbf{Why shift left by two?} Instructions are 4 bytes and word-aligned, so
\textbf{the low two bits of every instruction address are always zero}. Encoding
them would waste two bits of the immediate on information we already know.
Shifting instead gives \textbf{four times the branch range for free}.

\begin{quote}
That is a genuine architectural trade-off, and it fits in one line of
concatenation. It is also why the offset counts \textbf{instructions, not bytes}.
\end{quote}

\subsection*{2.4 Reuse: you already built the offset unit}

\code{ImmedXtnd} from Lab 4 already does sign-extend-then-shift. Set
\code{Xsel = '1'} and \code{Xlsft = '1'} and its output \emph{is} the offset ---
no new hardware. This is why \code{NIAG\_Rel} is one line:

\begin{lstlisting}
Rel <= std_logic_vector(unsigned(NSA) + unsigned(Xtnd32));
\end{lstlisting}

\textbf{And notice there is no subtractor.} Two's complement addition handles
backward branches by itself; a negative offset is just a large unsigned addend
that wraps.

\vspace{4pt}\hr

\section*{Example 3 --- Two ways to get it wrong}

All three of these compile, and all three produce plausible 32-bit addresses.

\begin{center}
\small
\begin{tabular}{l>{\raggedright\arraybackslash}p{6.4cm}>{\raggedright\arraybackslash}p{5.4cm}}
\toprule
version & formula & error shape \\
\midrule
\code{GOOD} & $(PC+4) + (\text{SignExt}(imm) \ll 2)$ & --- \\
\code{NOSHIFT} & $(PC+4) + \text{SignExt}(imm)$ & \textbf{off by a factor} of 4 \\
\code{BASEPC} & $PC + (\text{SignExt}(imm) \ll 2)$ & \textbf{off by a constant} 4 \\
\bottomrule
\end{tabular}
\end{center}

\subsection*{3.1 The symptoms, and which is worse}

\textbf{\code{NOSHIFT}} lands 2 bytes along instead of 8 --- \emph{in the middle
of an instruction}. The CPU then fetches a word straddling two instructions and
executes garbage. \textbf{Tell-tale: the target is not a multiple of 4.}

\textbf{\code{BASEPC}} is off by exactly 4, always, in the same direction. Every
branch lands one instruction early. \textbf{Tell-tale: the target is still
word-aligned and still nearly right}, which makes this the harder one to spot.
Loops run one instruction too many or too few.

\begin{quote}
\textbf{``Off by a factor'' and ``off by a constant'' are worth naming}, because
they are the two shapes almost every addressing bug takes --- here, in Lab 9's
memory addressing, and in Lab 12.
\end{quote}

\subsection*{3.2 The vector that separates them}

With $PC = \code{0x0C}$ and $imm = \code{0x0002}$:

\begin{center}
\begin{tabular}{ll}
\toprule
version & target \\
\midrule
\code{GOOD} & \code{0x18} \\
\code{NOSHIFT} & \code{0x12} \quad (not word-aligned) \\
\code{BASEPC} & \code{0x14} \quad (aligned, one instruction early) \\
\bottomrule
\end{tabular}
\end{center}

\textbf{One vector distinguishes all three}, provided the immediate is not zero
and the target is checked \emph{exactly} rather than ``looks about right''.

\textbf{Add a negative immediate too.} A zero-extension bug is invisible for
positive offsets, and every loop in the Lab 12 program ends with a backward
branch.

\vspace{4pt}\hr

\section*{Into the assignment}

\begin{center}
\begin{tabular}{lll}
\toprule
component & what it is & storage? \\
\midrule
\code{NIAG\_NSA} & $\code{NSA} = PC + 4$ & \textbf{no} --- combinational \\
\code{NIAG\_Rel} & $\code{Rel} = \code{NSA} + \code{Xtnd32}$ & \textbf{no} --- combinational \\
\code{NIAG\_Reg} & the program counter & \textbf{yes} --- the only register here \\
\bottomrule
\end{tabular}
\end{center}

\begin{quote}
\textbf{The program counter is the only storage element in this lab.} Both adders
must stay combinational: the CPU needs the next address \emph{during} the cycle in
which it fetches, because the result has to be ready at the edge that loads the
PC. Registering \code{NSA} makes the machine compute the next address from the
previous PC.
\end{quote}

Check in the Fitter report: with \code{NIAG\_Reg} as top level the register count
is \textbf{32}; with either adder as top level it is \textbf{0}.

\vspace{6pt}
\hr
{\small \textbf{Forward.} Lab 8 adds the condition logic and the multiplexer that
chooses among sequential, relative, absolute and return addresses. These three
components do not change --- they become three inputs of that mux.}

\end{document}
